\section{问题重述}

某小区微网由外部电网、光伏发电、用户负荷与一套电池储能系统组成。储能容量为12000 kWh，安全运行区间为1200--10800 kWh，最大充、放电功率均为5000 kW，2025年1月1日0:00储能量为6000 kWh。题目要求在满足负荷和储能约束的条件下，制定外购电与充放电计划，并依次考虑日前不确定性、日内光伏预报更新以及实时波动电价。

四问并非相互独立。问题一给出固定日的全部输入，重点是验证储能跨时段搬移电量的经济性；问题二要求每日0:00在信息不完全条件下做全年计划，并承担5倍电价的紧急购电；问题三允许在6:00、12:00和18:00利用新光伏预报调整尚未执行的计划；问题四进一步使外网电价随时间变化。故本文先建立公共物理模型，再逐层增加信息与结算结构。

\section{问题分析}

\subsection{问题一分析}

已知固定日的电价、负荷和光伏预测后，决策只在时间上耦合。低价时段或光伏富余时充电，高价时段放电，可以转移外购电时点；但若不设置日末库存约束，模型会无偿消耗初始电量。问题一适合用MILP表达能量平衡、库存递推、功率上限及充放电互斥，并用LP松弛作计算对照。

\subsection{问题二分析}

附件2给出的是事后实际曲线，不能在当天0:00直接作为已知量输入，否则高价紧急购电机制失去意义。本文用目标日以前的完整历史日形成点预测和联合残差场景；0:00计划在全部场景间共享，实际偏差由紧急购电与富余处置补救。评价时再用当天实际曲线回放，并将日末储能连续传至下一日。

\subsection{问题三分析}

日内新预报的价值来自信息变好，但计划调整又有不对称结算成本。应严格区分预报发布时间、目标时刻和预测提前量；已经执行的时段不可修改。本文以四个更新时刻重复求解剩余时域的两阶段随机模型，只执行到下一次更新，并用CVaR刻画紧急费用上尾风险。

\subsection{问题四分析}

波动电价使储能同时承担能量套利和预测风险缓冲。若未来电价在决策时不可知，应与净负荷共同构造条件场景；把全年实际电价提前输入只能作为完美信息下界。具体模型与结果占位见第9节。

\begin{figure}[H]
  \centering
  \begin{tikzpicture}[
  node distance=7mm and 7mm,
  box/.style={draw=black!65,rounded corners=2pt,align=center,minimum height=10mm,text width=27mm,fill=black!3,font=\small},
  hi/.style={box,fill=blue!7},
  arr/.style={-{Latex[length=2mm]},thick,draw=black!70}
]
\node[hi,text width=38mm] (physics) {公共物理层\\母线平衡、SOC递推\\容量、功率与互斥};
\node[box,right=of physics] (q1) {问题一\\确定性MILP\\固定日计划};
\node[box,right=of q1] (q2) {问题二\\历史预测与联合场景\\两阶段随机规划};
\node[box,right=of q2] (q3) {问题三\\四次预报更新\\随机MPC与CVaR};
\node[box,below=of q3] (q4) {问题四\\价格--净负荷联合场景\\波动电价调度};
\draw[arr] (physics) -- (q1);
\draw[arr] (q1) -- (q2);
\draw[arr] (q2) -- (q3);
\draw[arr] (q3) -- (q4);
\draw[arr] (physics.south) |- (q4.west);
\end{tikzpicture}

  \caption{四问共用物理层及信息递进关系}
  \label{fig:framework}
\end{figure}
